\chapter{บทนำ}

\section{ที่มาและความสำคัญ}
มหาวิทยาลัยเทคโนโลยีราชมงคลล้านนา มีภารกิจสนับสนุนการดำเนินงานโครงการหลวงและโครงการ
อันเนื่องมาจากพระราชดำริบนพื้นที่สูงภาคเหนือมาอย่างต่อเนื่อง เพื่อ\textbf{สืบสาน รักษา และต่อยอด}
พระราชปณิธานในการพัฒนาคุณภาพชีวิตของประชาชนบนพื้นที่สูง

ในอดีตการดำเนินงานลักษณะนี้กระจายอยู่ตามคณะ สถาบัน และวิทยาเขตต่าง ๆ ทำให้ยากต่อการ
ติดตามภาพรวม การรายงานผลต่อมหาวิทยาลัย และการวางแผนใช้ทรัพยากรร่วมกัน ในปีงบประมาณ 2569
มหาวิทยาลัยจึงรวมงานเหล่านี้เข้าเป็น \textbf{กลุ่มแผนงานใต้ร่มพระบารมี} ซึ่งเป็นแผนงานเดียว
ที่มีเป้าหมาย ตัวชี้วัด และระบบติดตามร่วมกัน

\section{วัตถุประสงค์ของแผนงาน}
\begin{enumerate}[leftmargin=1.6em,itemsep=2pt]
  \item วิจัยและพัฒนาเทคโนโลยีและนวัตกรรมที่ใช้ได้จริงในพื้นที่สูงและพื้นที่โครงการหลวง
  \item ยกระดับคุณภาพชีวิตและเสริมสร้างความเข้มแข็งของชุมชน ตามศาสตร์พระราชาและหลักปรัชญาเศรษฐกิจพอเพียง
  \item พัฒนากำลังคน สร้างอาชีพ และลดความเหลื่อมล้ำบนพื้นที่สูง
  \item สร้างกลไกการบริหารจัดการ การติดตามประเมินผล และเครือข่ายความร่วมมือที่ยั่งยืน
\end{enumerate}

\section{ขอบเขตของรายงาน}
รายงานฉบับนี้เป็น\textbf{รายงานความก้าวหน้ารอบ 10~เดือน} ครอบคลุมช่วงวันที่ 1 ตุลาคม 2568
ถึง 4 สิงหาคม 2569 คิดเป็นร้อยละ~85 ของระยะเวลาปีงบประมาณ

\begin{infobox}
\textbf{ข้อจำกัดของข้อมูลที่ต้องแจ้งให้ทราบ}\par
ตัวเลขตัวชี้วัดในรายงานฉบับนี้คือ\textbf{ค่าเป้าหมายที่แต่ละโครงการรับมาดำเนินการ}
เนื่องจากยังไม่มีการบันทึกผลสำเร็จรายกิจกรรมเข้าระบบติดตาม (0 จาก 268~กิจกรรม)
ผลสำเร็จที่ผ่านการตรวจรับจากหลักฐานจริงจะปรากฏในรายงานฉบับสมบูรณ์เมื่อสิ้นปีงบประมาณ
\end{infobox}

\section{วิธีการจัดทำรายงานและการตรวจสอบข้อมูล}
ข้อมูลในรายงานตรวจสอบตรงจากสามแหล่ง ได้แก่ ไฟล์งบประมาณฉบับ 4 สิงหาคม 2569
ไฟล์จัดสรรแผนงานและตัวชี้วัดฉบับปรับปรุง 20 กรกฎาคม 2569 และฐานข้อมูลระบบติดตามโครงการ

เพื่อป้องกันความคลาดเคลื่อนจากการคัดลอกด้วยมือ \textbf{ตารางรายโครงการในภาคผนวก ก และ ข
สร้างขึ้นโดยอัตโนมัติจากฐานข้อมูล} และยอดรวมทุกคอลัมน์ผ่านการตรวจสอบว่าบวกลงตัวตรงกับ
ยอดรวมในไฟล์งบประมาณต้นทาง
